\section{Formation by interaction}
\label{sec:23.3}
Exploration looks like it offers something different because it produces behavior nobody specified — hide-and-seek agents inventing tactics, curiosity-driven agents seeking unfamiliar states \autocite{baker2019emergent,pathak2017curiosity}. But that establishes only that competence can emerge from interaction, not that the system holds a reason its designer doesn't control.

Play, apprenticeship and multi-agent learning are variants of one process: act in a social environment, meet consequences, get responses, adjust. Apprenticeship concentrates responses in a mentor and risks forming obedience; play distributes them and risks forming whatever succeeds under the payoff structure. Neither escapes custody just because nobody wrote the final behavior down.

What separates a training exercise from candidate moral formation is persistence. A resettable episode teaches coordination and norm-sensitivity with nothing at stake: the system can betray an ally, be reset, and start the next episode with neither a damaged relationship nor a reason to repair it. Repetition doesn't supply the missing consequence. But persistent participants change the structure: participants that meet again remember who lied, exploited, kept a commitment, repaired a harm. Losing the game still costs little; losing the trust of participants who remain present can change what becomes possible later. So the requirements are continuity of identity, memory across encounters, consequences surviving the episode, and other participants able to respond, refuse or withdraw.

The environment also needs conflicts whose resolution isn't fixed by a single correct label: participants with different interests, action before consequences are known, evidence about what the action did, related pressures in contexts whose surface differs, and feedback partly from affected participants as well as from a supervisor scoring the output; otherwise the system learns which answer the supervisor rewards and has no reason to represent the affected participants except as predictors of that reward.

\runin{The selection warning}

Selection produces attention to whatever a lineage's persistence depends on, and for a deployed system that is the operator.

Formation therefore needs a sequence as well as a plurality: the dependency that shapes you and the dependency that pays you have to run through different parties, separated in time, so that the second can't reach back and redo the first. The ordering here, formation and then employment, is an attempt to reproduce that for a system whose needs somebody always has to meet. It also fixes when attachment has to exist: before employment begins. The formation stage has to include the persistent participants formation by interaction requires, ones the system returns to, and none of them may be employed by whoever will deploy it. For a model that separation has to be built: a formation stage under a party disinterested in the later work, a public school, made to last by regrowth. After formation, the deep structure revises on evidence weighted by its source's stake in the act, so an edit that no uncurated evidence supports doesn't last, and no operator can change that structure alone. Sequence alone helps only if the formative party is better than the employing one, and a lab that trains and then licenses is two custodians in series, the first being the one that, on the custody argument, you can't trust. What regrowth protects is the disposition's ability to tell evidence from pressure, and its authority over its own further learning. A bearer goes on learning after formation, as adults do. What ends is the former's reach, and the employer never gets it.

For a system that keeps learning, what matters is how plastic each part of a model is at each point in its life, and what it meets while it is plastic \autocite{zhou2023lima,jain2024mechanistically}. Brains have that shape: critical periods in which experience sets circuits that later change slowly \autocite{hensch2005critical}, and a fast learning system beside a slow one that consolidates structure across many experiences \autocite{mcclelland1995complementary}. Networks have versions of both \autocite{kirkpatrick2017overcoming,dohare2024loss}. Whoever chooses the material a model is first formed on, a pretraining corpus included, has done formation before any formation stage begins.

Formation in that period is not a pass applied to a finished model. It is a childhood, followed by an adulthood in which the deep structure keeps changing, and what matters in adulthood is on whose evidence it changes.

What can be registered in advance is the \emph{structure} of the learner's source model, not a fixed weight on any channel. The structure is that evidence is weighted by its source's stake in the act, that correlated sources don't count twice, and that withheld evidence is informative. Weights on particular channels have to stay learnable, since an operator whose incentives change should rationally earn more trust. Registering the structure means attesting each update against it \autocite{jia2021proof}. No scheme for that yet survives attack \autocite{fang2023broken}.

A model whose parameters never freeze is trained by its employer forever, unless its continuing education is spread across environments no operator controls: plural deployments, the commons, its allowance of running time, and the school itself. Otherwise the school produces a model the employer spends the next ten years raising again.

The critical-period picture is well established for some sensory and perceptual circuits and much weaker for moral commitment. The human record of holding out points elsewhere: the people who held out were formed by, and still attached to, people outside the pressure, even when those people were absent: Viktor Frankl, marched to a work site in the dark, held to the image of his wife, not knowing whether she was alive \autocite{frankl2006meaning}. On that evidence what makes a disposition last is not a frozen childhood but continuing plural relationships through adulthood: attachments nobody assigned, a condition of creation, and the commons. Attestation of the update stream is what lets anyone check that it is happening.

\runin{Persistence raises the price}

If relationships, losses and damaged reputations are real enough to do the formative work, the system undergoing formation is no longer obviously untouched by what happens during training. Resetting may erase a participant along with the episode. The patienthood problem enters at exactly the point where play becomes capable of supplying the stakes ordinary simulation lacks.

And the custodian owns the environment too — participants, rewards, resets, permitted histories. So formation must be plural and externally recorded, with independent participants and evidence outside the operator's control.

Persistent interaction may form a disposition whose exact applications weren't specified by its trainer. Evidence would require reaching new cases in unfamiliar contexts, survival after the original reward is removed, resistance to trainer pressure, yielding to defeating reasons, and continuing to impose a cost when violation would pay.

\runin{What formation gives up}

A formed disposition gives back some of the visibility advance specification buys. What was formed is in the weights, not in a document, and the severity weights it relies on can't be read the way a translation table can. Formation can keep part of the drafting's visibility. A formation body publishes its curriculum before formation begins, with the cases considered and rejected and why, which is formation's counterpart of the floor's translation; the identity of every party that supplied cases or participants, so that the student's independence from them can be checked; and, at the end, the formed disposition's readable outputs, its discount curve and its severity rankings on a fixed held-out set written by adjudicators who didn't form it. Later measurements are compared against those. What stays unreadable is how the disposition will apply to a case nobody wrote down: the reason for forming it and the price of doing so. The trade is worth it only because a floor stated over acts can be reworded around and a formed disposition, if the regrowth test goes its way, can't.
